\chapter{Short-Term Reversal}\label{ch:s1:short-term-reversal}

Buy yesterday's losers and sell yesterday's winners. On the synthetic market the residual version of this rule has a rank information coefficient of 0.041
against the next day's return, and the naive book built from it (gross exposure one, rebuilt at every close, in the 500 most liquid names) has a Sharpe ratio of
13.0 before costs. At ten basis points per unit traded it is $-3.4$. The strategy is real (it is one of the oldest anomalies in the literature, and the
synthetic market plants it) and nearly untradeable as written, because it turns over three quarters of its book every day. This chapter is about the one-day
reversal and its variants, and above all about how it is traded: with costs inside the optimiser it earns a Sharpe ratio of 1.7 after costs at \$100 million
and about nothing at \$1 billion. The build is \texttt{firm.reversal}.

\section{Why prices reverse: paying for immediacy}

\begin{definition}[Short-term reversal]\label{def:s1:short-term-reversal:reversal}
\emph{Short-term reversal}\index{short-term reversal} is the tendency of stocks with low returns over a short past window (a day, a week, a month) to
outperform, over the next such window, stocks with high returns; a short-term reversal strategy buys the recent losers and sells the recent winners.
\end{definition}

The effect has been in the literature for as long as there have been return files. Jegadeesh found significant negative first-order serial correlation in
monthly returns, the extreme deciles of a forecast built on it differing by 2.49\% a month over 1934--1987; Lehmann found that one week's winners and losers
reverse sizeably the next week, with profits that survived corrections for bid--ask spreads and plausible costs. Lehmann's reading has become the standard one:
the reversal is ``inefficiency in the market for liquidity around large price changes''. A trader who must sell now pushes the price down, the price comes
back as others absorb the flow, and whoever bought from the hurried seller is paid for providing immediacy. Nagel made the interpretation explicit, treating
reversal returns as a proxy for the returns from liquidity provision, and found them highest when liquidity is scarcest: expected returns and conditional
Sharpe ratios spike in financial turmoil, predictably with the VIX.

Two consequences shape everything that follows. First, the part of a price move that is news should not reverse, and the part that is pressure should, so a
good reversal signal tries to isolate pressure. Second, the reversal trader is a market maker at a daily horizon: the edge is the reward for trading, and the
strategy lives or dies by the cost of that trading.

\section{Daily reversal: raw, residual, industry-adjusted}

\begin{definition}[Residual reversal, industry-adjusted reversal]\label{def:s1:short-term-reversal:residual}
\emph{Residual reversal}\index{residual reversal} is \omterm{def:s1:short-term-reversal:reversal}{short-term reversal} computed on residual returns, the part of each stock's return that its factor
exposures do not explain. \emph{Industry-adjusted reversal}\index{industry-adjusted reversal} is reversal computed on each stock's return in excess of its
industry's average return.
\end{definition}

A stock's day is the market's day, its industry's day and its own. The market and industry parts are shared by hundreds of stocks and largely reflect news
about the economy; they are not what a hurried seller moves. The synthetic market (Book~7, chapter~5) plants a reversal of 5\% of each day's specific shock and
nothing on the common parts, so the cleaner the signal's estimate of the specific shock, the better it forecasts:

\begin{center}
\small
\begin{tabular}{@{}lcccc@{}}
\toprule
years 3 to 10 & raw & industry & residual & truth\\
\midrule
mean rank IC, all listed names & 0.037 & 0.038 & 0.041 & 0.044\\
mean rank IC, 500 most liquid names & 0.037 & 0.039 & 0.041 & 0.044\\
\midrule
naive book: Sharpe ratio before costs & 6.69 & 7.24 & 12.96 & \\
return a year; turnover a day & 27.0\%; 75\% & 28.2\%; 75\% & 30.3\%; 76\% & \\
Sharpe ratio after 10 bp & $-2.62$ & $-2.46$ & $-3.43$ & \\
break-even cost per unit traded & 7.2 bp & 7.5 bp & 7.9 bp & \\
\bottomrule
\end{tabular}
\end{center}

(The naive books trade the 500 most liquid names; ``industry'' is the industry-adjusted signal and ``truth'' the planted
expected return.) The information coefficient rises only from 0.037 to 0.041, but the residual book's Sharpe ratio almost doubles, because the raw book carries industry and
style exposures that add risk without return. Blitz, Huij, Lansdorp and Verbeek found the same on US data: conventional reversal strategies have dynamic
exposures to the Fama--French factors, and the residual version, without them, earns risk-adjusted returns twice as large, significant net of trading costs
even among large caps after 1990. Da, Liu and Schaumburg went further in the same direction, removing the part of the past return explained by cash-flow news
(measured with analysts' forecast revisions); their enhanced strategy's risk-adjusted return was four times the standard one's.

None of these gross numbers survives costs as a naive book. Every book in the table turns over about three quarters of its gross each day; at ten basis
points per unit of weight traded it pays 38\% a year against a gross return of 27--30\%. The break-even cost, the gross return divided by twice the turnover (Book~7, chapter~16),
is between 7 and 8 basis points for all three, and ten is already a low cost for a daily book (Book~7, chapter~16 made the same point with the raw signal).

\section{Intraday and overnight reversal}

\begin{definition}[Overnight reversal]\label{def:s1:short-term-reversal:overnight}
\emph{Overnight reversal}\index{overnight reversal} is the tendency of a stock's overnight return, from the previous close to the open, to be partly reversed
during the following trading session, and of its trading-session return to be partly reversed overnight; a strategy on it trades between the two parts of the
day.
\end{definition}

Splitting the day in two changes the picture. Lou, Polk and Skouras found strong continuation within each part (overnight returns predict overnight returns, and
intraday returns intraday returns, for years) together with an offsetting cross-period reversal, and in fourteen strategies that profits were earned either
entirely overnight or entirely intraday; for reversal, overnight. Akbas, Boehmer, Jiang and Koch counted, in each month, how often a positive overnight return
was followed by a negative trading-day reversal, and found that stocks with more such days (a more intense ``tug of war'' between overnight and daytime
clienteles) earned higher returns later. Inside the day, Heston, Korajczyk and Sadka traced the \omterm{def:s1:short-term-reversal:reversal}{short-term reversal} to temporary liquidity imbalances lasting
less than an hour and to bid--ask bounce, next to a striking continuation at intervals that are exact multiples of a day.

\texttt{firm.reversal} computes the overnight and intraday signals from open and close prices (functions \texttt{overnight} and \texttt{intraday});
the synthetic market has no open, so the chapter measures only the daily variants, and the two strategy files on them rest on the published record.

\section{Costs: where the strategy lives or dies}

A daily reversal book needs three things the naive book does not have. A forecast in units of return, so that it can be compared with costs: Grinold's rule
(Book~7, chapter~15) gives $\alpha_i = \mathrm{IC}\,\sigma_i\,z_i$, with $z_i$ the residual scaled by the stock's own specific volatility and standardised across
names, and the IC of 0.041 measured on year~2 only. A cost model: half a spread of 2 basis points plus square-root impact
$0.7\,\sigma_i\sqrt{Q_i/V_i}$ (Book~7, chapter~27). And an optimiser that compares the two for every name every day. With a diagonal risk model the problem
separates by name: each name's trade is a closed-form step that shrinks the desired trade by the linear cost and damps it by the impact, and neutrality to the
risk model's fifteen factors adds one small system for the multipliers (\cref{lst:s1:short-term-reversal:book}).

\omterm{def:s1:short-term-reversal:residual}{Residual reversal} on all names, earnings days skipped:

\begin{center}
\small
\begin{tabular}{@{}lccc@{}}
\toprule
fund size & \$100m & \$1bn & \$10bn\\
\midrule
costs ignored: Sharpe ratio before, after & 16.92, $-27.8$ & 16.92, $-31.7$ & 16.92, $-32.2$\\
costs ignored: trading costs a year & 246\% & 738\% & 2,293\%\\
costs inside: Sharpe ratio before costs & 8.42 & 3.97 & 1.72\\
costs inside: Sharpe ratio after costs & 1.72 & 0.09 & 0.19\\
costs inside: gross, costs, net (\% a year) & 10.23, 8.18, 2.05 & 2.44, 2.38, 0.05 & 0.45, 0.40, 0.05\\
costs inside: turnover a day; gross & 18.6\%; 0.91 & 3.9\%; 0.48 & 0.5\%; 0.18\\
net profit a year & \$2.0 million & \$0.5 million & \$4.8 million\\
\bottomrule
\end{tabular}
\end{center}

The cost-blind book trades 186\% of its capital a day at a gross exposure of 2.48: its gross Sharpe ratio of 16.9 is the best of any book in the chapter, and
at \$100 million it pays 246\% of its capital a year in costs. With costs inside, the book at \$100 million trades a tenth as much, keeps half its gross
Sharpe ratio, and nets 1.72. At \$1 billion it trades 3.9\% a day and earns its costs and nothing more; at \$10 billion it barely trades
(\cref{fig:s1:short-term-reversal:size}). The strategy's edge per trade is its information coefficient times a day's volatility, about sixteen basis
points for a name two standard deviations out, and a square-root cost reaches that on a trade of 1\% of the name's daily volume. The capacity of a one-day
reversal is therefore small: for this signal, somewhere between \$100 million and \$1 billion.

\begin{omfigure}
\begin{tikzpicture}
\begin{axis}[width=8.4cm,height=5.4cm,xlabel={\small fund size},ylabel={\small Sharpe ratio},tick label style={font=\footnotesize},
  xtick={8,8.477,9,9.477,10},xticklabels={\$100m,\$300m,\$1bn,\$3bn,\$10bn},ymin=-1,ymax=9,grid=major,grid style={gray!25},
  legend pos=north east,legend style={font=\scriptsize},table/col sep=comma]
\addplot[omDef,very thick,mark=*] table[x=log10_aum,y=sr_gross]{figdata/strategies-1/02-short-term-reversal/size.csv};
\addplot[omThm,very thick,mark=square*] table[x=log10_aum,y=sr_net]{figdata/strategies-1/02-short-term-reversal/size.csv};
\addplot[black!50,dashed] coordinates {(8,0) (10,0)};
\legend{before costs,after costs}
\end{axis}
\end{tikzpicture}

\omcaption{\omterm{def:s1:short-term-reversal:residual}{Residual reversal} with costs inside the optimiser, on the synthetic market: Sharpe ratio before and after costs against fund size. Data:
\texttt{s1\_reversal.traded}.}\label{fig:s1:short-term-reversal:size}
\end{omfigure}

One more filter made the difference between a small loss and a small profit. The synthetic market plants a post-earnings drift (Book~7, chapter~11), and a
stock that falls on bad earnings keeps drifting down: a reversal book that buys it is betting against news. Skipping names on their announcement day, which
is scheduled and public, costs almost nothing in information coefficient (0.040 against 0.041) and turns the book at \$1 billion from $-0.70\%$ a year
(Sharpe ratio $-1.04$) to $+0.05\%$ (\cref{fig:s1:short-term-reversal:earnings}). It is Da, Liu and Schaumburg's point on a small scale: moves explained by
fundamentals do not reverse.

\begin{omfigure}
\begin{tikzpicture}
\begin{axis}[width=9cm,height=5.4cm,xlabel={\small year},ylabel={\small cumulative net P\&L (\% of capital)},tick label style={font=\footnotesize},
  xmin=2,xmax=10,ymin=-7,ymax=4,grid=major,grid style={gray!25},legend pos=north east,legend style={font=\scriptsize},table/col sep=comma]
\addplot[omDef,very thick] table[x=year,y=filtered_pct]{figdata/strategies-1/02-short-term-reversal/earnings.csv};
\addplot[omThm,thick] table[x=year,y=unfiltered_pct]{figdata/strategies-1/02-short-term-reversal/earnings.csv};
\legend{announcement days skipped,every day}
\end{axis}
\end{tikzpicture}

\omcaption{The cost-aware reversal book at \$1 billion, net of costs, with and without the earnings filter. Data: \texttt{s1\_reversal.traded}.}\label{fig:s1:short-term-reversal:earnings}
\end{omfigure}

\section{The decline of reversal profits}

The public record agrees with the synthetic one on costs and adds a trend. Kenneth French's \omterm{def:s1:short-term-reversal:reversal}{short-term reversal} factor (value-weighted, low minus high
prior-month return, big and small stocks averaged, rebalanced monthly) earned 10.9\% a year from 1926 to 1989 with a Sharpe ratio of 0.93 ($t = 7.42$), and
1.6\% a year from 1990 to July 2026 with a Sharpe ratio of 0.13 ($t = 0.79$), before any trading costs; in the 2020s to date it has lost 5.2\% a year
(\cref{fig:s1:short-term-reversal:decline}). This is a monthly reversal on large, value-weighted portfolios, not the daily residual version, and it is not
evidence that the daily effect is gone; Nagel's results say that its returns concentrate in crises. But the long decline of the classic factor,
whatever its causes (publication, cheaper trading and more capital competing to provide liquidity are the usual candidates, and these data cannot
separate them), is the backdrop for every reversal book. The edge that remains belongs to
whoever trades it most cheaply.

\begin{omfigure}
\begin{tikzpicture}
\begin{axis}[width=10cm,height=5.0cm,xlabel={\small year},ylabel={\small 10-year mean (\% a year)},tick label style={font=\footnotesize},
  xmin=1935,xmax=2026,ymin=-5,ymax=30,grid=major,grid style={gray!25},table/col sep=comma,
  x tick label style={/pgf/number format/1000 sep={}}]
\addplot[omDef,very thick] table[x=year,y=mean_pct]{figdata/strategies-1/02-short-term-reversal/rolling.csv};
\addplot[black!50,dashed] coordinates {(1935,0) (2026,0)};
\addplot[omThm,thick,dashed] coordinates {(1990,-5) (1990,30)};
\node[font=\scriptsize,text=omThm,anchor=west] at (axis cs:1991,24) {published 1990};
\end{axis}
\end{tikzpicture}

\omcaption{Kenneth French's \omterm{def:s1:short-term-reversal:reversal}{short-term reversal} factor: the trailing 120-month mean return, annualised, at each December; the dashed line marks
Jegadeesh's and Lehmann's papers. Derived from the Kenneth R. French
Data Library (monthly ST\_Rev, 202607 CRSP file); the raw series is not redistributed.}\label{fig:s1:short-term-reversal:decline}
\end{omfigure}

\section{Strategy files}

\begin{strategyfile}{Daily residual reversal}\label{strat:s1:short-term-reversal:daily}
\sfield{Who pays you, and why} Traders who need to sell or buy now, and move prices to do it; the book is paid for absorbing their flow and waiting a day.
\sfield{Instruments and venues} Liquid single stocks, long and short, on their primary exchanges; the trades cross at the close or through the next session.
\sfield{Signal} Minus the day's residual return on a factor risk model, scaled by the stock's specific volatility; skip names with earnings or other scheduled
news that day.
\sfield{Sizing and execution} A forecast by Grinold's rule, a cost-aware optimiser neutral to the model's factors, trades worked into the close.
\sfield{Costs} The whole strategy: turnover of tens of per cent a day; half spread and impact on every trade; borrow on the shorts.
\sfield{How it dies} Costs rise or the book grows past its capacity; news days leak into the signal; other liquidity providers compete the edge away.
\sfield{Horizon, capacity, infrastructure} One to a few days; small capacity (for the chapter's signal, below \$1 billion); a daily risk model, a cost model
fitted on the firm's fills, a closing-auction execution system.
\sfield{Backtest honestly} Point-in-time residuals and universe; a cost model fitted on real fills, charged on every trade; the earnings filter known in advance;
results by size.
\sfield{Sources} Blitz, Huij, Lansdorp and Verbeek (2013); Da, Liu and Schaumburg (2014); this chapter's simulation.
\end{strategyfile}

\begin{strategyfile}{Intraday reversal to the close}\label{strat:s1:short-term-reversal:intraday}
\sfield{Who pays you, and why} Traders who push prices during the session with temporary imbalances; the book provides the other side and unwinds as the
imbalance clears.
\sfield{Instruments and venues} Liquid stocks and their lit and dark venues during the session.
\sfield{Signal} Minus the stock's return over the last minutes to an hour, relative to its industry or the market, with volume and imbalance conditions.
\sfield{Sizing and execution} Passive orders where possible; positions closed by the end of the day.
\sfield{Costs} Spread and adverse selection dominate; the edge is of the order of a spread.
\sfield{How it dies} Faster liquidity providers take the imbalance first; the move turns out to be news; bid--ask bounce mistaken for reversal.
\sfield{Horizon, capacity, infrastructure} Minutes to hours; small capacity; intraday data, an order-book simulator (Book~7, chapter~18) and low-latency
execution.
\sfield{Backtest honestly} Trade prices, not midquotes, in the signal, and midquotes in the evaluation; a queue-aware fill model.
\sfield{Sources} Heston, Korajczyk and Sadka (2010).
\end{strategyfile}

\begin{strategyfile}{Overnight reversal}\label{strat:s1:short-term-reversal:overnight}
\sfield{Who pays you, and why} The two clienteles of the tug of war: overnight and opening demand that pushes prices, and daytime traders who push them back.
\sfield{Instruments and venues} Liquid stocks; the opening and closing auctions.
\sfield{Signal} Minus the overnight return, traded over the session; or minus the session's return, traded overnight.
\sfield{Sizing and execution} Enter at the open auction and leave at the close, or the reverse; neutral to the market and industries.
\sfield{Costs} Two auction crossings a day; auction prices can move against large orders.
\sfield{How it dies} The continuation within each part of the day works against it; the pattern changes with who trades at the open.
\sfield{Horizon, capacity, infrastructure} Half a day; capacity limited by auction volume; open and close prices, auction imbalance data.
\sfield{Backtest honestly} Official open and close prices, not the first and last trades; auction costs; no signal from an open you could not trade at.
\sfield{Sources} Lou, Polk and Skouras (2019); Akbas, Boehmer, Jiang and Koch (2022).
\end{strategyfile}

\begin{strategyfile}{Reversal as a liquidity-provision overlay}\label{strat:s1:short-term-reversal:overlay}
\sfield{Who pays you, and why} The same hurried traders, but the reversal is not traded alone: it tilts the trades a slower book makes anyway.
\sfield{Instruments and venues} The firm's existing equity books.
\sfield{Signal} The reversal forecast added to the slower book's forecast in the optimiser, so that the book buys its planned purchases on down days and sells on
up days.
\sfield{Sizing and execution} Through the slower book's optimiser; the trades are the ones it would make anyway, with their timing shifted.
\sfield{Costs} Almost none of its own: the overlay changes when trades are made, not how many.
\sfield{How it dies} It does not die alone; it weakens when the slower book trades less.
\sfield{Horizon, capacity, infrastructure} Days; capacity set by the host book's turnover; the host book's optimiser.
\sfield{Backtest honestly} Compare the host book with and without the overlay at the same turnover and costs.
\sfield{Sources} Nagel (2012), on reversal returns as the return to liquidity provision; Book~7, chapter~27.
\end{strategyfile}

\section{Tutorial: seven before costs}\label{tut:s1:short-term-reversal:tut}

\begin{tutorial}
\textbf{Goal.} Measure the three daily reversals, then trade \omterm{def:s1:short-term-reversal:residual}{residual reversal} with costs inside the optimiser at three fund sizes. \textbf{End state:} the two
tables and \cref{fig:s1:short-term-reversal:size}.
\begin{tutsteps}
\item \textbf{Signals and ICs}: \texttt{s1\_reversal.signals}, \texttt{ic} for raw, industry-adjusted, residual and the truth, with and without earnings days.
\item \textbf{The naive books}: \texttt{naive(kind)} and \texttt{naive(kind, 0.001)}, and \texttt{breakeven}.
\item \textbf{The cost-aware book.} A closed-form step per name and Newton's method for the factor-neutrality multipliers.
\omcode{code/firm/reversal/firm_reversal.py}{69}{97}{The day's cost-aware, factor-neutral book.}\label{lst:s1:short-term-reversal:book}
\item \textbf{Every day}: forecast, optimise, charge the cost, drift the book.
\omcode{code/strategies-1/02-short-term-reversal/python/s1_reversal.py}{137}{151}{The daily loop of the traded book.}\label{lst:s1:short-term-reversal:loop}
\item \textbf{Run} \texttt{traded(aum, costs, 1.0, skip\_earnings)} at \$100 million, \$1 billion and \$10 billion, and \texttt{fig\_reversal.py}.
\end{tutsteps}
\textbf{What to change next.} Smooth the forecast (Book~7, chapter~27) and see whether a slower book beats a faster one; restrict the universe to the 500 most
liquid; multiply the cost in the optimiser by two, as if every trade had to be unwound, and explain why it does not help here.
\end{tutorial}

\section{Build: the reversal family}\label{bld:s1:short-term-reversal:reversal}

\begin{build}
\bfield{Purpose} The signals of the reversal family and a daily book that trades them only where they beat their costs.
\bfield{Interface} \texttt{raw(r)}, \texttt{industry\_adjusted(r, industry, listed)}, \texttt{residual(r, X, w, listed)}, \texttt{overnight(open\_,
prev\_close)}, \texttt{intraday(close, open\_)}, \texttt{forecast(z, ic, sigma)}, \texttt{book(alpha, var, w0, gamma, aum, sigma, adv, half\_spread, eta,
costs, round\_trip, X)}.
\bfield{Rules} Signals from the day's data only; forecasts in return units; the cost in the optimiser is the cost charged; neutrality to the exposures held
exactly.
\bfield{Acceptance tests} \texttt{code/firm/reversal/tests/}: the signals by hand; the residual orthogonal to the exposures; without costs, the book equals the
closed-form mean--variance solution; costs shrink the trades; factor neutrality holds to $10^{-10}$ with costs.
\bfield{Stretch} A full factor covariance in the optimiser (Book~7, chapter~27's \texttt{cost\_aware}); multi-period trading toward an aim portfolio (Book~7,
chapter~26); the overnight and intraday signals on real open and close prices.
\end{build}

\begin{omsources}
\item N.~Jegadeesh, ``Evidence of predictable behavior of security returns'', \emph{Journal of Finance} 45(3), 1990.
\item B.~N.~Lehmann, ``Fads, martingales, and market efficiency'', \emph{Quarterly Journal of Economics} 105(1), 1990.
\item S.~Nagel, ``Evaporating liquidity'', \emph{Review of Financial Studies} 25(7), 2012.
\item D.~Blitz, J.~Huij, S.~Lansdorp and M.~Verbeek, ``Short-term residual reversal'', \emph{Journal of Financial Markets} 16(3), 2013.
\item Z.~Da, Q.~Liu and E.~Schaumburg, ``A closer look at the short-term return reversal'', \emph{Management Science} 60(3), 2014.
\item S.~L.~Heston, R.~A.~Korajczyk and R.~Sadka, ``Intraday patterns in the cross-section of stock returns'', \emph{Journal of Finance} 65(4), 2010.
\item D.~Lou, C.~Polk and S.~Skouras, ``A tug of war: overnight versus intraday expected returns'', \emph{Journal of Financial Economics} 134(1), 2019.
\item F.~Akbas, E.~Boehmer, C.~Jiang and P.~D.~Koch, ``Overnight returns, daytime reversals, and future stock returns'', \emph{Journal of Financial Economics}
145(3), 2022.
\item Kenneth R. French Data Library, Short-Term Reversal Factor (ST Rev), monthly.
\end{omsources}

\section{Exercises}

\begin{exercise}[$\star$]\label{exo:s1:short-term-reversal:1}
The residual naive book returns 30.3\% a year before costs and turns over 76\% of its gross a day. Compute its break-even cost per unit traded.
\end{exercise}

\begin{exercise}[$\star$]\label{exo:s1:short-term-reversal:2}
A stock's midquote does not move for two days; the spread is 10 basis points. The last trade of day one is at the bid and of day two at the ask. What return
does a close-to-close file record for day two, and what does a reversal signal on trade prices conclude?
\end{exercise}

\begin{exercise}[$\star$]\label{exo:s1:short-term-reversal:3}
With square-root impact $0.7\,\sigma\sqrt{Q/V}$, what is the impact of a \$10 million trade in a stock with daily volatility 2\% and \$25 million of daily volume?
\end{exercise}

\begin{exercise}[$\star\star$]\label{exo:s1:short-term-reversal:4}
The IC is 0.04 and a stock's daily volatility 2\%. What is Grinold's forecast for a name two standard deviations out, and how large a trade in the stock of
exercise~3 can it pay for, with a half spread of 2 basis points?
\end{exercise}

\begin{exercise}[$\star\star$]\label{exo:s1:short-term-reversal:5}
Why does the \omterm{def:s1:short-term-reversal:residual}{residual reversal} book have almost twice the raw book's Sharpe ratio when its information coefficient is only 10\% higher?
\end{exercise}

\begin{exercise}[$\star\star$]\label{exo:s1:short-term-reversal:6}
Explain why a reversal book that buys stocks on the day of a bad earnings announcement loses money in a market with post-earnings drift, and why skipping
announcement days is not look-ahead.
\end{exercise}

\begin{exercise}[$\star\star\star$]\label{exo:s1:short-term-reversal:7}
\emph{Coding.} Run \texttt{traded(3e8, True, 1.0, True)}. Report the net Sharpe ratio, and explain why the net Sharpe ratio is not a monotone function of fund
size in \cref{fig:s1:short-term-reversal:size}.
\end{exercise}

\begin{exercise}[$\star\star\star$]\label{exo:s1:short-term-reversal:8}
\emph{Find the flaw.} ``Our \omterm{def:s1:short-term-reversal:residual}{residual reversal} backtest has a Sharpe ratio of 13. Even if live trading halves it, a \$1 billion allocation is safe.''
\end{exercise}

\section{Problem: Seven Before Costs}

\begin{problem}[{Weekend problem --- a strategy that lives in its costs}]\label{pb:s1:short-term-reversal:1}
The chapter's reversal signals on \texttt{firm.synthmkt}, and the reversal factor on real data.

\medskip\noindent\textbf{Part I --- The effect.}
\begin{enumerate}
\item Define \omterm{def:s1:short-term-reversal:reversal}{short-term reversal} and give the liquidity-provision explanation.
\item What did Jegadeesh and Lehmann find?
\item What did Nagel show about when reversal pays?
\item Why should news not reverse?
\end{enumerate}

\medskip\noindent\textbf{Part II --- The signals.}
\begin{enumerate}[resume]
\item Give the ICs of raw, industry-adjusted and \omterm{def:s1:short-term-reversal:residual}{residual reversal}, and the truth's.
\item Give the naive books' Sharpe ratios before and after ten basis points.
\item Why does the residual book do so much better before costs?
\item What did Blitz and co-authors and Da, Liu and Schaumburg find?
\end{enumerate}

\medskip\noindent\textbf{Part III --- Trading it.}
\begin{enumerate}[resume]
\item How is the forecast put in return units?
\item Describe the optimiser and why it separates by name.
\item What does the cost-blind book trade and pay?
\item What does the earnings filter do, and what is it worth at \$1 billion?
\end{enumerate}

\medskip\noindent\textbf{Part IV --- The verdict.}
\begin{enumerate}[resume]
\item State the \emph{named result}: the net Sharpe ratio of \omterm{def:s1:short-term-reversal:residual}{residual reversal} at \$100 million and \$1 billion with costs inside the optimiser, against the
naive book's.
\item Where is the strategy's capacity, and why so small?
\item What happened to French's reversal factor after 1989?
\item Why is that not proof the daily effect is gone?
\item What do overnight and intraday decompositions add?
\item Which strategy file would you run with \$50 million, and which with \$5 billion?
\item What is the single most important input to a reversal backtest?
\item In one sentence: what is \omterm{def:s1:short-term-reversal:reversal}{short-term reversal}, as a business?
\end{enumerate}
\end{problem}

\section{Interview questions}

\begin{interviewq}[$\star$ \iqroles{researcher, trader}]\label{iq:s1:short-term-reversal:1}
Why do short-term returns reverse, and who pays the reversal trader?
\end{interviewq}

\begin{interviewq}[$\star\star$ \iqroles{researcher}]\label{iq:s1:short-term-reversal:2}
Residual versus raw reversal: what is the difference, and why does it matter?
\end{interviewq}

\begin{interviewq}[$\star\star$ \iqroles{trader}]\label{iq:s1:short-term-reversal:3}
A reversal strategy has a gross Sharpe ratio of 7 and a break-even cost of 15 basis points. Would you trade it? What would you ask?
\end{interviewq}

\begin{interviewq}[$\star\star$ \iqroles{developer, researcher}]\label{iq:s1:short-term-reversal:4}
How would you build a daily cost-aware optimiser for a thousand names that runs in seconds?
\end{interviewq}

\begin{interviewq}[$\star\star$ \iqroles{risk}]\label{iq:s1:short-term-reversal:5}
When does a reversal book lose money fast, and what would you monitor?
\end{interviewq}

\begin{interviewq}[$\star\star\star$ \iqroles{researcher}]\label{iq:s1:short-term-reversal:6}
Under square-root impact, derive the trade size at which a one-day forecast $\alpha$ just pays for its cost, for a name with volatility $\sigma$, daily volume
$V$ and half spread $h$.
\end{interviewq}
